\selectlanguage{ngerman}%
\frenchspacing%
\begin{ledgroupsized}[r]{120mm}%
\footnotesize%
\pstart%
\noindent%
\textbf{Überlieferung:}%
\pend%
\end{ledgroupsized}%
\begin{ledgroupsized}[r]{114mm}%
\footnotesize%
\pstart%
\parindent -6mm%
\makebox[6mm][l]{\textit{L}}%
Aufzeichnung:
LH XXXV~6,\;23 Bl.\;3.
Ein Blatt 8\textsuperscript{o};
Ränder zum Teil beschnitten;
römisches\indexO{Rom} Papier %
(ohne Wasserzeichen).
Eine halbe Seite auf Bl.\;\href{\Scan{LH_35_06_23_003r}}{3\,r\textsuperscript{o}}
und zwei Drittel einer Seite auf Bl.\;\href{\Scan{LH_35_06_23_003v}}{3\,v\textsuperscript{o}};
Textfolge nicht eindeutig. %
Bl.\;\href{\Scan{LH_35_06_23_003r}}{3\,r\textsuperscript{o}} überliefert ferner im unteren Seitenteil den in diesem Band veröffentlichten Text N.\;\ref{58656-B}, %
im oberen Seitenteil folgendes Textfragment von Leibnizens Hand:
\textit{%
\sout{Datum est quod exposito}
Datus est angulus qui exposito aequalis vel idem est.
Ita Ratio data est, quae eadem rationi magnitudin}%
\bracks{\textit{um}}
(Handschrift: \textit{magnitudini})
\textit{datarum.
Data directio, quae est eadem directioni expositae.}%
\pend%
\end{ledgroupsized}%
\vspace{5mm}%
\begin{ledgroup}%
\footnotesize%
\pstart%
\noindent%
\textbf{Datierungsgründe:}
Die Datierung der Aufzeichnung N.\;\ref{58656-A} beruht auf Eigenschaften des Textträgers sowie auf inhaltlichen Betrachtungen.
Siehe die ausführliche Begründung auf S.\;\pageref{58656-A_dat}.%
\pend%
\end{ledgroup}%
\selectlanguage{latin}%
\frenchspacing%
\vspace{8mm}%
\normalsize%
\pstart%
\noindent%
\edlabel{LH_35_06_23_003_003r_anfg-1}%
\edtext{}{%
\xxref{LH_35_06_23_003_003r_anfg-1}{LH_35_06_23_003_003r_anfg-2}%
\lemma{\lbrack\href{\Scan{LH_35_06_23_003r}}{3\,r\textsuperscript{o}}\rbrack}%
\Bfootnote{%
\textit{(1)}~Eadem posita mole \lbrack/\rbrack\ %
\textit{(2)}~\textlangle Vires intra datum spatium co\textrangle\ \lbrack/\rbrack\ %
\textit{(3)}~Vires %
\textit{(4)}~Pondera intra datum spatium %
\textit{(a)}~coercentes %
\textit{(b)}~coercentia sunt %
\bars{ut \textit{erg.}} densitates seu moles, %
\textit{(aa)}~densitat %
\textit{(bb)}~seu iisdem positis %
\textit{(aaa)}~ponderibus %
\textit{(bbb)}~spatiis pondera %
\lbrack...\rbrack\ spatiorum et densitatum. \lbrack/\rbrack\ %
\textit{(aaaa)}~Pondera eadem mole coercentia sunt ut \lbrack/\rbrack\ %
\textit{(bbbb)}~Pondera eandem molem %
\bars{ejusdem caloris \textit{erg.}} coercentia sunt ut %
\textit{(aaaaa)}~spatia seu reci %
\textit{(bbbbb)}~densitates seu reciproce ut spatia.%
~\textit{L}%
}}%
\lbrack%
\href{\Scan{LH_35_06_23_003r}}%
{3\,r\textsuperscript{o}}%
\rbrack\
\lbrack%
\textit{Nachfolgend kleingedruckter Text in L gestrichen:}%
\rbrack\
\pend%
\vspace{0.5em}%
\footnotesize%
\pstart%
\noindent%
Pondera intra datum spatium coercentia%
\indexS{pondus coercens molem}%
\indexS{spatium datum}
sunt ut densitates seu moles,%
\indexS{densitas molis coercitae}%
\indexS{moles coercita pondere}
seu iisdem positis spatiis
pondera sunt ut densitates.
\pend%
\pstart%
Iisdem positis ponderibus
spatia sunt reciproce ut densitates.
Ergo densitates sunt in ratione composita
ex ponderum directa et spatiorum reciproca.
\pend%
\pstart%
Seu pondera sunt in ratione composita
spatiorum et densitatum.
\pend%
\vspace{1.0em}%
\normalsize%
\pstart%
\noindent%
Pondera eandem molem ejusdem caloris coercentia%
\indexS{pondus coercens molem}%
\indexS{moles coercita pondere}%
\indexS{moles ejusdem caloris}%
\indexS{calor molis coercitae}
sunt ut densitates%
\indexS{densitas molis coercitae}
seu reciproce ut spatia.%
\indexS{spatium molis coercitae}%
\edlabel{LH_35_06_23_003_003r_anfg-2}
\pend%
\pstart%
Pondera eandem molem intra idem spatium coercentia%
\indexS{pondus coercens molem}%
\indexS{moles coercita pondere}
sunt ut calores%
\indexS{calor molis coercitae}
(calores enim sunt ut auctae%
\indexS{calor ut densitas aucta}%
\indexS{densitas molis aucta}
\edtext{%
densitates;
at%
}{%
\lemma{densitates;}%
\Bfootnote{%
\textit{(1)}~et %
\textit{(2)}~at%
~\textit{L}%
}}
pondera materiam%
\indexS{pondus coercens materiam}%
\indexS{materia coercita pondere}
\edtext{%
duplo densiorem%
}{%
\lemma{duplo}%
\Bfootnote{%
\textit{(1)}~densa\textlangle m\textrangle\ %
\textit{(2)}~densiorem%
~\textit{L}%
}}
intra idem spatium coercentia%
\indexS{spatium materiae coercitae}%
\indexS{densitas materiae coercitae}
\edlabel{LH_35_06_23_003_saltus-1}%
\edtext{}{%
\xxref{LH_35_06_23_003_saltus-1}{LH_35_06_23_003_saltus-2}%
\lemma{sunt}%
\Bfootnote{%
\textit{(1)}~ut densitates %
\textit{(2)}~dupla). \lbrack/\rbrack\ %
\textit{(a)}~Pondera \lbrack/\rbrack\ %
\textit{(b)}~Itaque pondera \lbrack\href{\Scan{LH_35_06_23_003v}}{3\,v\textsuperscript{o}}\rbrack\ %
\textit{(c)}~Vires quae dupl \lbrack/\rbrack\ %
\textit{(d)}~Pondera materiam %
\lbrack...\rbrack\ sunt dupla.%
~\textit{L}%
}}%
sunt dupla).
\lbrack%
\href{\Scan{LH_35_06_23_003v}}%
{3\,v\textsuperscript{o}}%
\rbrack%
\pend%
\pstart%
Pondera materiam duplo densiorem%
\indexS{pondus coercens materiam}%
\indexS{materia coercita pondere}
intra idem spatium coercentia%
\indexS{spatium materiae coercitae}%
\indexS{densitas materiae coercitae}
sunt dupla.%
\edlabel{LH_35_06_23_003_saltus-2}
Nam pondera
\edtext{%
duplam%
}{%
\lemma{duplam}%
\Bfootnote{%
\textit{erg.\;L}%
}}
materiam aeque densam
intra duplum spatium coercentia
sunt eadem
\edtext{%
per \dots\
Ergo pondera
duplam materiam aeque densam
intra simplum%
\indexS{pondus coercens materiam}%
\indexS{materia coercita pondere}%
\indexS{spatium materiae coercitae}%
\indexS{densitas materiae coercitae}
\edtext{%
spatium coercentia,
sunt dupla
per \dots%
}{%
\lemma{spatium}%
\Bfootnote{%
\textit{(1)}~coercentia, id est materiam %
\textit{(2)}~coercentia, %
\textit{(a)}~id est %
\textit{(b)}~sunt dupla per \dots%
~\textit{L}%
}}%
\
Sed hoc ipsum est
materiam duplo densiorem
intra idem spatium coercere%
\indexS{materia coercita pondere}%
\indexS{spatium materiae coercitae}%
\indexS{densitas materiae coercitae}
per def.\;\dots%
\indexS{definitio}%
\
Ergo
\edlabel{LH_35_06_23_003_003v_schluss-1}%
\edtext{}{%
\xxref{LH_35_06_23_003_003v_schluss-1}{LH_35_06_23_003_003v_schluss-2}%
\lemma{duplum}%
\Bfootnote{%
\lbrack/\rbrack\ %
\textit{(1)}~Aequalibus positis molibus, vires coercentes sunt ut densitates %
\lbrack/\rbrack\ %
\textit{(2)}~Vires quae moles aequal %
\textit{(3)}~Similes sunt aequales, hoc est%
\textit{(a)}~Si mole %
\textit{(b)}~Si spatia %
\lbrack...\rbrack\ coercentes sunt ut %
\textit{(aa)}~spatia seu ut densitates,
\textit{(bb)}~densitates, seu reciproce ut spatia.%
~\textit{L}%
}}
duplum%
}{%
\lemma{per \bracks{...} duplum}%
\Cfootnote{%
Siehe zu den Andeutungen
\textsc{G.\,W.\;Leibniz},
\textit{Dynamica}, pars\;II, sectio\;I, prop.\;22--25
(%
\textit{CFP}\;I, S.\;327--331~[A], 570--572~[B]; II.2, S.\;395--400~[C]).\cite{01354}%
}}
\lbrack%
\textit{Text bricht ab.}%
\rbrack\
\pend%
\begin{edfigure}[1.5em][1.0em]
\includegraphics%
{gesamttex/edit_VIII-4/images/LH_35_06_23_003_d1_003v.pdf}
\edcaption[0.25em]
{\lbrack \textit{Fig.\:1}\rbrack}%
\end{edfigure}%
\pstart%
Similes sunt aequales,
hoc est%
\lbrack:\rbrack\
Si spatia sint densitatibus reciproce proportionalia,
vires coercentes sunt ut densitates,
seu reciproce ut spatia.%
\indexS{vis coercens materiam}%
\indexS{spatium materiae coercitae}%
\indexS{densitas materiae coercitae}%
\edlabel{LH_35_06_23_003_003v_schluss-2}
\pend%
\pstart%
Si densitates sint eaedem
et spatia sint utcunque,
vires sunt%
\indexS{vis coercens materiam}%
\indexS{spatium materiae coercitae}%
\indexS{densitas materiae coercitae}
\edlabel{LH_35_06_23_003_003v_math-1}%
eaedem.
\edtext{}{%
\xxref{LH_35_06_23_003_003v_math-1}{LH_35_06_23_003_003v_math-2}%
\lemma{eaedem.}%
\Bfootnote{%
\textit{(1)}~Si vires sint eaedem %
\textit{(2)}~Si spatia sint eadem %
\textit{(3)}~Ergo pondera %
\lbrack...\rbrack\ aut moles. %
\textit{(a)}~$HL = CD$ %
\textit{(b)}~\textit{DB} dupl.\;\textit{FB} %
\lbrack...\rbrack\ $DB = \text{bis}\;LB$ %
\textit{(aa)}~rursusque \textlangle\textit{DB}\textrangle\ %
\textit{(bb)}~et contra dens.\;\textit{LM} $=$ bis dens.\;\textit{DB}.%
\textit{(aaa)}~Ergo %
\textit{(aaaa)}~$CD~=$ %
\textit{(bbbb)}~$HL = \text{bis}\;CD$. %
\textit{(cccc)}~$GF = \text{bis}\;CD$. %
\textit{(bbb)}~Ergo $HL = \text{bis}\;CD$.%
~\textit{L}%
}}%
Ergo pondera sunt ut densitates absolutae,
quaecunque sint spatia aut moles.%
\indexS{pondus coercens molem}%
\indexS{moles coercita pondere}%
\indexS{spatium molis coercitae}%
\indexS{densitas molis coercitae}
\pend%
\pstart%
\textit{DB} dupl.\;\textit{FB}
et dens.\;\textit{DB} $=$ dens.\;\textit{FB}.
Ergo
$CD = EF$
et
$DB = \text{bis}\;LB$
et contra
dens.\;\textit{LM} $=$ bis dens.\;\textit{DB}.
Ergo
$HL = \text{bis}\;CD$.
Ergo $HL = \text{bis}\;EF$.%
\edlabel{LH_35_06_23_003_003v_math-2}%
\pend%
\begin{edfigure}[1.750em][0em]
\includegraphics%
{gesamttex/edit_VIII-4/images/LH_35_06_23_003_d2_003v.pdf}
\edcaption[-0.25em]{\lbrack \textit{Fig.\:2}\rbrack}%
\end{edfigure}%
\newpage%
\pstart%
\noindent%
\centering%
*     *     *
\pend%
\vspace{2.0em}%
\selectlanguage{german}%
\frenchspacing%
\footnotesize%
\pstart%
\noindent%
\textls{Ausführliche Datierungsgründe:}%
\pagelabel{58656-A_dat}
Die titellose und undatierte Aufzeichnung N.\;\ref{58656-A} umfasst flüchtige Bemerkungen zum Verhältnis von Druck, Dichte und Volumen bei einer komprimierbaren Masse, die einige Verwandschaft mit Leibnizens \textit{Dynamica} aufweisen
(vgl. bes. pars II, sectio I, prop.\;22--25:\cite{01354}
\textit{CFP}\;I, S.\;327--331~[A], 570--572~[B]; II.1, S.\;395--400~[C]).\cite{01043} %
Dies legt die Annahme nahe, dass N.\;\ref{58656-A} zu einem Zeitpunkt ab dem Sommer 1689 entstand, als Leibniz zwischen Rom\indexO{Rom} und Florenz\indexO{Florenz} am Entwurf der \textit{Dynamica} arbeitete.
\pend%
\pstart%
Eine Bestätigung dieser Annahme und zugleich einen Grund dafür, die Entstehungszeit der Aufzeichnung N.\;\ref{58656-A} weiter einzugrenzen, liefert ihr Träger.
Denn obwohl dieser kein Wasserzeichen aufweist, lässt sich seine römische\indexO{Rom} Herkunft anhand sichtbarer Eigenschaften des Papiers (wie Farbe und Anordnung der Stege) feststellen:
Auf gleichem Papier sind etwa auch die in diesem Band veröffentlichten Aufzeichnungen N.\;\ref{58258}, N.\;\ref{58218} und N.\;\ref{58235} %
überliefert (bei N.\;\ref{58258} und N.\;\ref{58235} liegt auch ein inhaltlicher Zusammenhang mit Leibnizens \textit{Dynamica} vor; siehe die ausführlichen Datierungsgründe auf S.\;\pageref{58258_dat} u.\;\pageref{58235_dat}).
Ein Vergleich mit weiteren Handschriften aus dem Leibniz-Nachlass
-- etwa mit %
den Auszügen aus Giovanni Alfonso Borellis\indexN{\textls{Borelli} (Borellus), Giovanni Alfonso 1608\textendash1679} \textit{De vi percussionis}\cite{01001} (\textit{LSB} VIII,\;3 N.\;73),\cite{02602}
mit den Aufzeichnungen \textit{De motu elaterii se restituentis} (\textit{LSB} VIII,\;3 N.\;28\textsubscript{1} u.\;28\textsubscript{2})\cite{02603}
und mit %
dem Konzept des Briefes an Christoph Gudenus\indexN{\textls{Gudenus}, Johann Christoph 1632\textendash1705} vom 6.\;August 1689 (\textit{LSB} I,\;5 N.\;250)\cite{01297}
-- zeigt,
dass N.\;\ref{58656-A} sowie N.\;\ref{58258}, N.\;\ref{58218} und N.\;\ref{58235} auf Papier verfasst wurden, das Leibniz während seines Aufenthaltes in Rom\indexO{Rom} benutzte (das Wasserzeichen weist eine umkreiste sechsstrahlige Sonne mit Kreuz und Buchstaben F auf). %
Auf gleichem Papier sind auch etliche Bogen der handschriftlichen A-Fassung der \textit{Dynamica}\cite{01354} überliefert.
Somit ist eine Entstehung sämtlicher Aufzeichnungen N.\;\ref{58656-A}, N.\;\ref{58258}, N.\;\ref{58218} und N.\;\ref{58235} zu der Zeit, als Leibniz in Rom\indexO{Rom} an der ersten Fassung der \textit{Dynamica}\cite{01354} arbeitete, sehr wahrscheinlich.
\pend%
\pstart%
Die vorliegende Aufzeichnung N.\;\ref{58656-A} ist demgemäß auf den Zeitraum zwischen Sommer und November 1689 zu datieren (siehe zur Entstehung von Leibnizens \textit{Dynamica}\cite{01354} \textsc{Robinet} 1988, S.\;81--83 u.\;259-262).\cite{01240}
Eine frühere (ab Leibnizens Ankunft in Rom\indexO{Rom} im Mai 1689) oder eine spätere Abfassungszeit (höchstens bis Ende März 1690, als Leibniz Italien\indexO{Italien} verlassen hat) sind nicht ausgeschlossen, aber unwahrscheinlich.%
\pend%
\selectlanguage{latin}%
\frenchspacing%
\normalsize%
\newpage%
